\documentclass[12pt,a4paper]{article}
\usepackage[utf8]{inputenc}
\usepackage{geometry}
\geometry{a4paper, margin=1in}
\usepackage{hyperref}
\usepackage{booktabs}
\usepackage{longtable}
\usepackage{xcolor}
\usepackage{listings}

\hypersetup{
    colorlinks=true,
    linkcolor=blue,
    urlcolor=cyan,
    pdftitle={Phase 1 Status Report - DAS 839}
}

\lstset{
    basicstyle=\ttfamily\small,
    breaklines=true,
    backgroundcolor=\color{gray!10},
    frame=single,
    rulecolor=\color{gray!30}
}

\title{\textbf{Phase 1 Status Report}\\DAS 839 -- NoSQL Systems\\Multi-Pipeline ETL and Reporting Framework}
\author{Team Submission}
\date{\today}

\begin{document}
\maketitle

\section{Current Verified Status}
A fresh end-to-end Hive run was executed successfully using:
\begin{lstlisting}
python3 controller.py --pipeline hive --batch-size 100000
\end{lstlisting}
Latest verified run id: \texttt{16}.

Persisted metadata for run 16 (most recent):
\begin{itemize}
        \item pipeline: hive
    \item batch size: 500000
    \item total records: 3461613
    \item malformed records: 8
    \item number of batches: 7
    \item average batch size: 494516.14
    \item runtime: 91095 ms
\end{itemize}
We verified correctness by checking that aggregated request counts match total valid records and that batch IDs remain contiguous across the run.
Although some sample batches can show zero errors, the global aggregation for the full run confirms non-zero error counts.

\section{Clarification: Local Mode vs HDFS}
The current Phase 1 Hive prototype uses a hybrid single-machine setup:
\begin{itemize}
    \item \textbf{Input storage/source}: HDFS path \texttt{hdfs:///nasalogs/raw/}
    \item \textbf{Compute execution mode}: Hive-on-MapReduce local mode (single machine)
    \item \textbf{Output files}: local filesystem under \texttt{/tmp/etl\_output/run\_\{run\_id\}}
\end{itemize}

So data is read from HDFS, but computation is performed in local MR mode on this PC.

\section{Where Output Files Exist}
For run 13, output files exist at:
\begin{lstlisting}
/tmp/etl_output/run_13/
  malformed/000000_0
  q1/000000_0
  q2/000000_0
  q3/000000_0
  step1.log
  step2.log
\end{lstlisting}
Aggregated output is also loaded into PostgreSQL tables:
\texttt{run\_metadata}, \texttt{q1\_daily\_traffic}, \texttt{q2\_top\_resources}, \texttt{q3\_hourly\_errors}.

\section{Implemented vs Planned Pipeline Scope}
\begin{itemize}
    \item \textbf{Implemented in Phase 1}: Hive pipeline end-to-end (extract/transform/batch/query/load/report).
    \item \textbf{Planned for Phase 2}: MapReduce, Pig, and MongoDB backends with the same ETL semantics and reporting schema.
    \item \textbf{Controller readiness}: CLI selection supports all four pipeline choices already; non-Hive paths are currently placeholders.
\end{itemize}
The interface contract supports four pipeline choices, but only Hive is implemented at present.

\section{Hive ETL Workflow (What Happens and Why)}
\subsection{Extract}
\begin{itemize}
    \item Hive external table \texttt{nasa\_raw(line STRING)} reads raw NASA log lines from HDFS.
    \item No manual pre-cleaning/pre-conversion is done outside the pipeline.
\end{itemize}

\subsection{Transform}
\begin{itemize}
    \item Parse fields using \texttt{REGEXP\_EXTRACT}: host, date components, method, path, protocol, status, bytes.
    \item Convert \texttt{DD/Mon/YYYY} into ISO \texttt{YYYY-MM-DD}.
    \item Derive \texttt{log\_hour} as integer.
    \item Normalize bytes: \texttt{'-'}, empty, or null to \texttt{0}.
    \item Records failing core parse conditions are excluded from aggregates and counted as malformed.
\end{itemize}

\subsection{Batching}
\begin{itemize}
    \item Assign batch id with:
\begin{lstlisting}[language=SQL]
CEIL(ROW_NUMBER() OVER (ORDER BY line) / ${BATCH_SIZE})
\end{lstlisting}
    \item This guarantees sequential batch IDs starting from 1 and includes a final partial batch.
\end{itemize}

\subsection{Load}
\begin{itemize}
    \item Hive writes malformed, Q1, Q2, Q3 outputs as TSV files.
    \item \texttt{result\_loader.py} inserts outputs into PostgreSQL reporting tables.
    \item \texttt{controller.py} persists final runtime after DB write completion.
\end{itemize}

\subsection{Report}
\begin{itemize}
    \item \texttt{reporter.py} reads PostgreSQL and displays metadata + query outputs.
\end{itemize}

\section{Mandatory Query Implementation and Output Schemas}
\subsection{Q1: Daily Traffic Summary}
Output schema used:
\begin{itemize}
    \item \textbf{batch\_id}, log\_date, status\_code, request\_count, total\_bytes
\end{itemize}

\subsection{Q2: Top Requested Resources}
Output schema used:
\begin{itemize}
    \item \textbf{batch\_id} (set to 0 by design), resource\_path, request\_count, total\_bytes, distinct\_host\_count
\end{itemize}
This query is computed globally and not per batch.

\subsection{Q3: Hourly Error Analysis}
Output schema used:
\begin{itemize}
    \item \textbf{batch\_id}, log\_date, log\_hour, error\_request\_count, total\_request\_count, error\_rate, distinct\_error\_hosts
\end{itemize}
The current report stores \texttt{error\_rate} as a percentage in the range 0--100, rounded to 2 decimals.
Execution time is represented in \texttt{run\_metadata.runtime\_ms} and tied to query results through \texttt{run\_id}.

\section{Assignment Compliance Checklist (Phase 1)}
\begin{longtable}{p{0.10\textwidth}p{0.17\textwidth}p{0.65\textwidth}}
\toprule
\textbf{Item} & \textbf{Status} & \textbf{Evidence / Notes} \\
\midrule
1.1 & Met (Hive path) & End-to-end ETL + PostgreSQL load + reporting demonstrated via controller for Hive. Same dataset and logical query definitions maintained. \\
1.2 & Met (Hive path) & Required fields parsed; bytes fallback to 0 for '-', empty, null; malformed records counted (run 13 malformed=8). \\
1.3 & Partially met (Phase 1 scope) & CLI supports all four pipeline choices; Hive implemented and working; other backends are placeholders for now (to be implemented in Phase 2). \\
1.4 & Met (Hive path) & Q1/Q2/Q3 implemented with required output columns and validated output rows. \\
1.5 & Met (Hive path) & Batch IDs sequential; average batch size reported from metadata formula; runtime ends after DB write completion. \\
1.6 & Met (architecture), Partial (full implementation) & Coherent tool architecture is integrated; backend switching design in place; only Hive fully implemented in Phase 1 prototype. \\
1.7 & Substantially met & Phase 1 minimum prototype objective is substantially met via one working pipeline, with the second pipeline planned. Hive prototype and architecture documentation are ready. \\
1.8 & Process requirement & Strict LMS deadline requirement; submit status report and Git link before Tuesday 3:30 PM. \\
\bottomrule
\end{longtable}

\section{Pipeline Selection Interface (Phase 1 Behavior)}
The CLI already accepts exactly one pipeline argument:
\begin{lstlisting}
--pipeline hive | mapreduce | pig | mongodb
\end{lstlisting}
Current behavior:
\begin{itemize}
    \item \textbf{hive}: fully working
    \item \textbf{mapreduce, pig, mongodb}: explicit placeholders in Phase 1 (clean message + non-zero exit)
\end{itemize}

\section{Fair and Comparable Cross-Pipeline Plan}
To ensure fair comparison in Phase 2:
\begin{itemize}
    \item same source NASA dataset
    \item same canonical parse and cleaning semantics
    \item same Q1/Q2/Q3 logic and output schemas
    \item same batch size per experiment
    \item same runtime boundary definition (read start to DB write completion)
    \item same reporting schema in PostgreSQL
    \item same query contract: Q1 and Q3 batched, Q2 global top-20 with \texttt{batch\_id=0}
\end{itemize}

\section{Workflow Diagrams}
\subsection{Overall ETL + Reporting Flow}
\begin{verbatim}
[NASA logs source]
      |
      v
[controller.py --pipeline X --batch-size N]
      |
      +--> [Hive | MapReduce | Pig | MongoDB]
                 |
                 v
      [Parsed + Batched + Q1/Q2/Q3 outputs]
                 |
                 v
          [result_loader.py]
                 |
                 v
            [PostgreSQL]
                 |
                 v
             [reporter.py]
\end{verbatim}

\subsection{Detailed Hive ETL Workflow}
\begin{verbatim}
Extract:
    HDFS raw files -> Hive external table nasa_raw(line)

Transform:
    REGEXP_EXTRACT for host/date/hour/method/path/protocol/status/bytes
    bytes cleanup: '-', empty, null -> 0
    malformed lines counted separately

Batch:
    batch_id = CEIL(ROW_NUMBER() OVER (ORDER BY line) / batch_size)

Aggregate:
    Q1 (daily status summary)
    Q2 (top 20 resources)
    Q3 (hourly error analysis)

Load:
    TSV outputs -> result_loader.py -> PostgreSQL

Report:
    reporter.py reads PostgreSQL and prints metadata + query results
\end{verbatim}

\subsection{Hive Parsing + Batching Strategy}
\begin{verbatim}
nasa_raw(line)
   |
   +-- REGEXP_EXTRACT(host, date, hour, method, path, protocol, status, bytes)
   |
   +-- bytes cleanup: '-', empty, null -> 0
   |
   +-- malformed count tracked separately
   |
   +-- batch_id = CEIL(ROW_NUMBER() OVER (ORDER BY line) / N)
   |
   v
nasa_parsed(host, log_date, log_hour, ..., bytes, batch_id)
\end{verbatim}

\subsection{Hive File Sequence with Inputs and Outputs}
\begin{verbatim}
controller.py
    input : --pipeline hive, --batch-size N, DB config
    output: run_id created in run_metadata
        |
        v
pipelines/hive/run_hive.sh
    input : N, run_id, output base dir
    output: /tmp/etl_output/run_<run_id>/{malformed,q1,q2,q3,step1.log,step2.log}
        |
        +--> 01_create_tables.hql
        |      input : HDFS raw log lines
        |      output: nasa_parsed table
        |
        +--> 02_queries.hql
                     input : nasa_raw + nasa_parsed
                     output: malformed/q1/q2/q3 TSV files
        |
        v
common/result_loader.py
    input : TSV output folders
    output: inserts into PostgreSQL result tables
        |
        v
controller.py runtime finalization
    output: updates run_metadata.runtime_ms
        |
        v
common/reporter.py
    output: terminal display of metadata + query samples
\end{verbatim}

\subsection{Cross-Pipeline Fairness Design (Phase 2 Plan)}
\begin{verbatim}
For each pipeline P in {Hive, MapReduce, Pig, MongoDB}:
    same dataset source
    same parse rules
    same cleaning rules
    same Q1/Q2/Q3 definitions
    same batch size per experiment
    same runtime boundary
    same PostgreSQL output schema

Then compare outputs across run_ids for equivalence.
\end{verbatim}

\section{Phase 1 Conclusion and Next Steps}
\section{CSV Export Feature (Phase 1 Enhancement)}
All query results are automatically exported to CSV files after each run. This provides easy spreadsheet access to results without requiring database queries.

\subsection{CSV Output Structure}
Results are organized in \texttt{output/hive/run\_\{run\_id\}/} with standardized filenames:
\begin{itemize}
    \item \texttt{00\_run\_metadata.csv} -- pipeline name, batch size, record counts, runtime, etc.
    \item \texttt{01\_q1\_daily\_traffic.csv} -- all 1681 rows for run 16 (daily traffic by batch/date/status)
    \item \texttt{02\_q2\_top\_resources.csv} -- all 20 top resources (global aggregation)
    \item \texttt{03\_q3\_hourly\_errors.csv} -- all 9285 rows for run 16 (hourly errors by batch/date/hour)
\end{itemize}

\subsection{Implementation}
The controller automatically calls \texttt{common/export\_results.py} after PostgreSQL load completes. This module:
\begin{enumerate}
    \item Queries PostgreSQL for all results with the current run\_id.
    \item Writes CSV files with headers using Python's csv module.
    \item Creates run-specific folders for persistence and traceability.
\end{enumerate}

\subsection{Access Methods}
Results can be accessed via:
\begin{lstlisting}
# View in terminal
cat output/hive/run_16/01_q1_daily_traffic.csv | head -20

# Open in spreadsheet
libreoffice output/hive/run_16/02_q2_top_resources.csv

# Count rows
wc -l output/hive/run_16/03_q3_hourly_errors.csv
\end{lstlisting}

Every run automatically generates these files, eliminating the need for manual export or scripting.

\section{Phase 1 Conclusion and Next Steps}
Phase 1 minimum prototype objective is substantially met via one working pipeline with complete output formats, with the second pipeline planned.
Phase 1 minimum prototype objective is substantially met via one working pipeline, with the second pipeline planned.

Next steps for Phase 2:
\begin{enumerate}
    \item Implement MapReduce backend using the same semantics and output schema.
    \item Implement Pig backend.
    \item Implement MongoDB backend.
    \item Add cross-run equivalence checks across backends.
    \item Produce comparative analysis (correctness, runtime, implementation complexity).
\end{enumerate}

\end{document}
